\begin{table}[t]
\centering
\caption{Playable protagonist-form variables used in the coding framework. The categories are used to test whether archetypal distributions differ by how protagonism is enacted.}
\label{tab:protagonist_form_variables}
\small
\setlength{\tabcolsep}{4pt}
\renewcommand{\arraystretch}{1.08}
\begin{tabularx}{\textwidth}{p{2.2cm}p{3.1cm}X}
\toprule
\textbf{Variable} & \textbf{Category} & \textbf{Operational cue and examples} \\
\midrule
Type & Fixed character & Authored named protagonist with relatively stable identity, e.g. Lara Croft, Kratos, Aloy, Arthur Morgan. \\
Type & Silent/projective avatar & Minimally voiced or weakly biographical avatar whose role is partly produced through player projection, e.g. Gordon Freeman, Chell, Link. \\
Type & Customizable protagonist & Player-configured identity within a designed role, e.g. Shepard, V, Tav, Dragonborn, Tarnished. \\
Type & Implicit/systemic role & No conventional protagonist; the player acts as ruler, manager, creator, planner, or institutional role, e.g. \game{Civilization}, \game{SimCity}, \game{Crusader Kings}. \\
\midrule
Structure & Single / dual / ensemble & Protagonism is concentrated in one character, split between two, or distributed across several playable figures. \\
Structure & Party or avatar-plus-companions & A central avatar interacts with a companion field that shapes the narrative and archetypal position. \\
Structure & Systemic position & Protagonism is enacted through rules, resources, interfaces, and consequences rather than through a character body. \\
\midrule
Arc & Player-shaped arc & Degree to which choices, endings, moral systems, or role-playing decisions can redirect the protagonist's archetypal emphasis. \\
\bottomrule
\end{tabularx}
\end{table}
